\subsection{Convergence Theorems}

% **Role:** modern integration and almost-everywhere analysis.
%
% Monotone convergence, Fatou, dominated convergence, uniform integrability preview.

This subsection collects the main limit-interchange tools: monotone convergence, Fatou, and dominated convergence.

\begin{theorem} \label{thm:lebesgue-product-main-source-todo}\label{thm:lebesgue-convergence-theorems}\label{thm:monotone-convergence} \hlt{Monotone Convergence}\\
Let $E\in\mathcal{M}$, let $f_n,f$ be measurable on $E$, and assume $0\leq f_n\uparrow f$. Then $\int_E f_n\,d\mu\uparrow\int_E f\,d\mu$.
\begin{proof}
By monotonicity, $\int_Ef_n\,d\mu\uparrow L\leq\int_Ef\,d\mu$.\\
For $s\in\mathcal{S}^+(f)$ and $0<c<1$, set $E_n\defeq\{x\in E:f_n(x)\geq cs(x)\}$.\\
Since $f_n\uparrow f$ and $s\leq f$, $E_n\uparrow E$, and $\int_Ef_n\,d\mu\geq\int_{E_n}f_n\,d\mu\geq c\int_{E_n}s\,d\mu$.\\
Continuity from below for $\phi_s(A)\defeq\int_As\,d\mu$ gives $L\geq c\int_Es\,d\mu$. Let $c\uparrow1$, then take $\sup_{s\in\mathcal{S}^+(f)}$.\\
\hlt{Structural Commentary}: Nonnegative integration commutes with monotone limits.\\
\hlt{Key Insight}: Simple minorants of the limit are eventually captured below the sequence.
\end{proof}
\end{theorem}

\begin{theorem} \label{thm:nonnegative-series-integral} \hlt{Series Form of Monotone Convergence}\\
Let $E\in\mathcal{M}$, let $f_n\geq0$ be measurable, and set $f\defeq\sum\limits_{n=1}^{\infty}f_n$ on $E$. Then $\int_E f\,d\mu=\sum\limits_{n=1}^{\infty}\int_E f_n\,d\mu$.
\begin{proof}
Set $S_N\defeq\sum\limits_{n=1}^{N}f_n$. Then $0\leq S_N\uparrow f$.\\
By Theorem \ref{thm:monotone-convergence}, $\int_Ef\,d\mu=\lim\limits_{N\to\infty}\int_ES_N\,d\mu =\lim\limits_{N\to\infty}\sum\limits_{n=1}^{N}\int_Ef_n\,d\mu =\sum\limits_{n=1}^{\infty}\int_Ef_n\,d\mu$.\\
\hlt{Structural Commentary}: Nonnegative series are monotone partial sums.\\
\hlt{Key Insight}: Integrate partial sums, then pass upward.
\end{proof}
\end{theorem}

\begin{theorem} \label{thm:fatou-lemma} \hlt{Fatou's Lemma}\\
Let $E\in\mathcal{M}$ and let $f_n\geq0$ be measurable. Then $\int_E\liminf_{n\to\infty}f_n\,d\mu\leq\liminf_{n\to\infty}\int_Ef_n\,d\mu$
\begin{proof}
Set $g_n\defeq\inf_{i\geq n}f_i$. Then $0\leq g_n\uparrow\liminf_i f_i$ and, for every $i\geq n$, $g_n\leq f_i$.\\
Hence $\int_Eg_n\,d\mu\leq\inf_{i\geq n}\int_Ef_i\,d\mu$. By Theorem \ref{thm:monotone-convergence}, let $n\to\infty$ to obtain the claim.\\
\hlt{Structural Commentary}: Fatou compares the liminf to tail lower envelopes.\\
\hlt{Key Insight}: Apply monotone convergence to $g_n=\inf_{i\geq n}f_i$.
\end{proof}
\end{theorem}

\begin{theorem} \label{thm:dominated-convergence-lebesgue} \hlt{Dominated Convergence}\\
Let $E\in\mathcal{M}$, $f_n$ measurable, $f_n\to f$ a.e. on $E$, and $\abs{f_n}\leq g\in\mathscr{L}^1(\mu)$ on $E$.\\
Then $f\in\mathscr{L}^1(\mu)$ on $E$ and $\lim_{n\to\infty}\int_Ef_n\,d\mu=\int_Ef\,d\mu$.
\begin{proof}
Choose $N\subseteq E$ null such that $f_n\to f$ on $E\setminus N$; replacing all functions by their restrictions off $N$ does not change integrals. Thus assume $f_n\to f$ everywhere. Then $\abs{f}\leq g$, so $f\in\mathscr{L}^1(\mu)$. Put $a_n\defeq\int_Ef_n\,d\mu$.\\
For $g+f_n\geq0$, by Fatou, $\int_E(g+f)\,d\mu\leq\liminf_n\int_E(g+f_n)\,d\mu =\int_Eg\,d\mu+\liminf_na_n$; $\int_Ef\,d\mu\leq\liminf_na_n$.\\
For $g-f_n\geq0$, by Fautou, $\int_E(g-f)\,d\mu\leq\liminf_n\int_E(g-f_n)\,d\mu=\int_Eg\,d\mu-\limsup_na_n$; $\limsup_na_n\leq\int_Ef\,d\mu$. Thus $\liminf_na_n=\limsup_na_n=\int_Ef\,d\mu$.\\
\hlt{Structural Commentary}: Domination turns signed convergence into two Fatou inequalities.\\
\hlt{Key Insight}: Apply Fatou to $g+f_n$ and $g-f_n$ to trap liminf and limsup.
\end{proof}
\end{theorem}

\begin{corollary} \label{cor:monotone-convergence-ae} \hlt{Monotone Convergence Almost Everywhere}\\
Let $E\in\mathcal{M}$ and $0\leq f_n\uparrow f$ a.e. on $E$. Then $\int_Ef_n\,d\mu\uparrow\int_Ef\,d\mu$.
\begin{proof}
Choose $N\subseteq E$ with $\mu(N)=0$ and $f_n\uparrow f$ on $E\setminus N$. Set $g_n\defeq f_n\mathbf{1}_{E\setminus N}$ and $g\defeq f\mathbf{1}_{E\setminus N}$.\\
Then $0\leq g_n\uparrow g$, so $\int_Eg_n\,d\mu\uparrow\int_Eg\,d\mu$. Since $f_n=g_n$ and $f=g$ off $N$, the integrals over $N$ vanish.\\
\hlt{Structural Commentary}: Monotone convergence ignores null exceptional sets.\\
\hlt{Key Insight}: Remove the exceptional set, prove the theorem there, then add back zero.
\end{proof}
\end{corollary}

\begin{example} \label{ex:counting-measure-integral-is-series} \hlt{Counting Measure Integrals Are Series}\\
Let $\mu$ be counting measure on $(\N,2^\N)$. For $f:\N\to[0,+\infty]$, set $s_N\defeq\sum_{n=1}^{N}f(n)\mathbf{1}_{\{n\}}$; then $0\leq s_N\uparrow f$, with the case $f(k)=+\infty$ giving both sides $+\infty$. Hence $\int_{\N} f\,d\mu=\lim_N\int_{\N}s_N\,d\mu =\lim_N\sum\limits_{n=1}^{N}f(n) =\sum\limits_{n=1}^{\infty}f(n)$.\\
\hlt{Structural Commentary}: Counting measure turns integration into summation.\\
\hlt{Key Insight}: Finite-support truncations are the simple-function partial sums.
\end{example}

\begin{example} \label{ex:monotone-fatou-boundary-tests} \hlt{Boundary Tests for Monotone Convergence and Fatou}\\
On $([0,+\infty),\mathcal{B},m)$:
\begin{enumerate}[label=(\roman*)]
\setlength{\itemsep}{0pt}
\item $f_n\defeq\mathbf{1}_{[0,n]}\uparrow\mathbf{1}_{[0,+\infty)}$ and $\int f_n\,dm=n\uparrow+\infty$; monotone convergence allows $+\infty$.
\item $g_n\defeq n^{-1}\mathbf{1}_{[n,+\infty)}\downarrow0$ uniformly, but $\int g_n\,dm=+\infty$ $\forall n$; decreasing convergence needs finite mass.
\item $w_n\defeq-n^{-1}\mathbf{1}_{[0,n]}\rightrightarrows0$, but $\int w_n\,dm=-1$; Fatou can fail without $f_n\geq0$ or $f_n\geq-h\in L^1$.
\end{enumerate}
\end{example}

\begin{theorem} \label{thm:fatou-with-integrable-lower-bound} \hlt{Fatou with an Integrable Lower Bound}\\
Let $h\geq0$, $h\in\mathscr{L}^1(\mu)$, and let $f_n$ be measurable with $f_n\geq-h$.\\\
Then, whenever the signed integrals are defined in $\overline{\R}$, $\int \liminf\limits_{n\to\infty} f_n\,d\mu\leq\liminf\limits_{n\to\infty}\int f_n\,d\mu$.
\begin{proof}
Set $F_n\defeq f_n+h\geq0$. Fatou gives $\int \liminf\limits_nF_n\,d\mu\leq\liminf\limits_n\int F_n\,d\mu$.\\
Since $\liminf_nF_n=h+\liminf_nf_n$ and $\int F_n\,d\mu=\int f_n\,d\mu+\int h\,d\mu$, subtract the finite term $\int h\,d\mu$.\\
\hlt{Structural Commentary}: Fatou survives controlled negative parts.\\
\hlt{Key Insight}: Add an integrable floor to make the sequence nonnegative.
\end{proof}
\end{theorem}

\begin{remark} \label{rem:negative-fatou-example-no-integrable-floor} \hlt{Why the Negative Example Has No Floor}\\
For $w_n\defeq-n^{-1}\mathbf{1}_{[0,n]}$, any $h\in L^1$ with $-h\leq w_n$ $\forall n$ must satisfy $h\geq n^{-1}$ on $[0,n]$.\\
Thus $h\geq k^{-1}$ on $[k-1,k]$, and $\int_0^N h\,dm\geq H_N\defeq\sum\limits_{k=1}^{N}\frac1k\to+\infty$, contradicting $h\in L^1$.
\end{remark}

\begin{perspective} \label{per:domination-as-limit-integral-safeguard} \hlt{Domination Is the Safeguard}\\
$f_n\to f$ a.e. need not imply $\int f_n\to\int f$; domination $\abs{f_n}\leq g\in L^1$ restores the limit interchange.
\end{perspective}

\begin{example} \label{ex:dominated-convergence-improper-integral} \hlt{A Nonuniform Limit with Vanishing Integrals}\\
For $x>0$, set $f_n(x)\defeq e^{-nx}x^{-1/2}$. Then $f_n\to0$ pointwise, but $\sup_{0<x<\delta}f_n(x)=+\infty$ for every $\delta>0$.\\
For $n\geq1$, $0\leq f_n\leq g\defeq e^{-x}x^{-1/2}$; $g\leq x^{-1/2}\mathbf{1}_{(0,1]}+e^{-x}\mathbf{1}_{[1,\infty)}\in\mathscr{L}^1(0,\infty)$.\\
Thus dominated convergence gives $I_n\defeq\int_0^\infty e^{-nx}x^{-1/2}\,dx\to0$.\\
\hlt{Structural Commentary}: Uniform convergence is unnecessary once one integrable envelope controls the sequence.\\
\hlt{Key Insight}: Domination handles the singular behavior at $0$ and the tail at infinity simultaneously.
\end{example}

\begin{corollary} \label{cor:dominated-parameter-continuity} \hlt{Continuity Under a Lebesgue Integral}\\
Let $f_t(x)\defeq f(x,t)$ be measurable for each $t\in[a,b]$. Fix $t_0\in[a,b]$.\\
If $f_t\to f_{t_0}$ a.e. as $t\to t_0$ and $\abs{f_t}\leq g\in\mathscr{L}^1(\mu)$ for all $t$, then $\lim\limits_{t\to t_0}\int_X f_t\,d\mu=\int_X f_{t_0}\,d\mu$.\\
If $t\mapsto f(x,t)$ is continuous on $[a,b]$ for a.e. $x$, then $F(t)\defeq\int_X f_t\,d\mu$ is continuous on $[a,b]$.
\begin{proof}
For $t_n\to t_0$, set $u_n\defeq f_{t_n}$.\\
Then $u_n\to f_{t_0}$ a.e. and $\abs{u_n}\leq g$, so dominated convergence gives $\int_Xu_n\,d\mu\to\int_Xf_{t_0}\,d\mu$.\\
The sequential criterion gives the first claim; applying it for every $t_0$ gives continuity of $F$.\\
\hlt{Structural Commentary}: Parameter continuity reduces to dominated convergence along sequences.\\
\hlt{Key Insight}: The parameter disappears once $t_n$ is frozen into a sequence of functions.
\end{proof}
\end{corollary}

\begin{theorem} \label{thm:differentiation-under-lebesgue-integral} \hlt{Differentiation Under a Lebesgue Integral}\\
Let $f_t(x)\defeq f(x,t)$ be measurable for each $t\in[a,b]$. Suppose $f_{t_0}\in\mathscr{L}^1(\mu)$ for some $t_0\in[a,b]$, $\partial_tf$ exists on $X\times(a,b)$,\\
and $\abs{\partial_tf(x,t)}\leq g(x)$ for some $g\in\mathscr{L}^1(\mu)$ and all $(x,t)\in X\times(a,b)$.\\
Then $f_t\in\mathscr{L}^1(\mu)$ for every $t\in[a,b]$, and, with $F(t)\defeq\int_Xf_t\,d\mu$, $F'(t)=\int_X \partial_tf(x,t)\,d\mu$, $t\in(a,b)$.
\begin{proof}
For fixed $t$, $x\mapsto\partial_tf(x,t)$ is measurable as the pointwise limit of measurable difference quotients.\\
By the mean value theorem, $\abs{f_t-f_{t_0}}\leq \abs{t-t_0}g$, so $\abs{f_t}\leq\abs{f_{t_0}}+\abs{t-t_0}g\in\mathscr{L}^1(\mu)$.\\
If $t_n\to t\in(a,b)$ and $t_n\neq t$, then $\frac{F(t_n)-F(t)}{t_n-t} =\int_X\frac{f_{t_n}-f_t}{t_n-t}\,d\mu$.\\
Integrand tends pointwise to $\partial_tf(\cdot,t)$ and has modulus $\leq g$; dominated convergence gives derivative formula.\\
\hlt{Structural Commentary}: Bounded parameter derivatives dominate the difference quotients.\\
\hlt{Key Insight}: Differentiating the integral is just applying dominated convergence to secant slopes.
\end{proof}
\end{theorem}

\begin{corollary} \label{cor:single-point-differentiation-under-lebesgue-integral} \hlt{A Pointwise Difference-Quotient Test}\\
Let $f_t(x)\defeq f(x,t)$ be measurable for each $t\in[a,b]$, let $f_{t_0}\in\mathscr{L}^1(\mu)$ for some $t_0\in[a,b]$, and fix $t_1\in[a,b]$.\\
Suppose $\partial_tf(x,t_1)$ exists for every $x$ and $\exists g\in\mathscr{L}^1(\mu)$ such that, for $t\neq t_1$, $Q_t(x)\defeq\frac{f_t(x)-f_{t_1}(x)}{t-t_1}$, $\abs{Q_t}\leq g$.\\
Then $F(t)\defeq\int_Xf_t\,d\mu$ is differentiable at $t_1$ and $F'(t_1)=\int_X\partial_tf(x,t_1)\,d\mu$.
\begin{proof}
If $t_0\neq t_1$, then $\abs{f_{t_1}}\leq\abs{f_{t_0}}+\abs{t_0-t_1}g$; otherwise this is given.\\
Hence $f_{t_1}\in\mathscr{L}^1$, and $\abs{f_t}\leq\abs{f_{t_1}}+\abs{t-t_1}g$ gives $f_t\in\mathscr{L}^1$. For any $t_n\to t_1$, $t_n\neq t_1$, $Q_{t_n}\to\partial_tf(\cdot,t_1)$ pointwise and $\abs{Q_{t_n}}\leq g$, so dominated convergence gives $[F(t_n)-F(t_1)]/(t_n-t_1)\to\int_X\partial_tf(x,t_1)\,d\mu$.\\
\hlt{Structural Commentary}: One parameter point needs only domination of the secants into that point.\\
\hlt{Key Insight}: A global derivative bound is sufficient, but not necessary, for a one-point formula.
\end{proof}
\end{corollary}

\begin{corollary} \label{cor:riemann-parameter-integral-interchange} \hlt{One Lebesgue and One Riemann Integral}\\
Let $f_t(x)\defeq f(x,t)$ be measurable for each $t\in[a,b]$, let $t\mapsto f(x,t)$ be continuous on $[a,b]$ for a.e. $x$, and suppose $\abs{f_t}\leq g\in\mathscr{L}^1(\mu)$ $\forall t$. Then, with Riemann integrals in $t$, $\int_a^b\left[\int_X f_t\,d\mu\right]dt =\int_X\left[\int_a^b f(x,t)\,dt\right]d\mu(x)$.
\begin{proof}
Set $h_t(x)\defeq\int_a^t f(x,s)\,ds$. Then $h_t$ is measurable, $\abs{h_t}\leq(b-a)g$, and $\partial_th_t=f_t$.\\
By differentiation under the integral, $H(t)\defeq\int_Xh_t\,d\mu$ satisfies $H'(t)=\int_Xf_t\,d\mu$.\\
Hence $\int_a^b\left[\int_Xf_t\,d\mu\right]dt=H(b)-H(a)=\int_X\left[\int_a^b f(x,t)\,dt\right]d\mu(x)$.
\hlt{Structural Commentary}: Finite parameter integrals move through by differentiating primitives.\\
\hlt{Key Insight}: Integrate first in the parameter up to $t$, then differentiate the resulting Lebesgue integral.
\end{proof}
\end{corollary}

\begin{corollary} \label{cor:improper-riemann-parameter-integral-interchange} \hlt{Improper Parameter Interchange}\\
Let $f_t(x)\defeq f(x,t)$ be measurable for each $t\in\R$, and let $t\mapsto f(x,t)$ be continuous for a.e. $x$.\\
Suppose $\exists g,h\in\mathscr{L}^1(\mu)$ such that $\abs{f_t}\leq g$ $\forall t$ and $\int_\R\abs{f(x,t)}\,dt\leq h(x)$ for a.e. $x$.\\
Put $G(t)\defeq\int_Xf_t\,d\mu$; $H(x)\defeq\int_\R f(x,t)\,dt$. Then $H\in\mathscr{L}^1(\mu)$; $\int_\R G(t)\,dt=\int_XH\,d\mu$, $\int_\R G\defeq\lim\limits_{R\to+\infty}\int_{-R}^{R}G(t)\,dt$.
\begin{proof}
For $R>0$, finite-interval interchange gives $\int_{-R}^{R}G(t)\,dt=\int_XH_R\,d\mu$ and $H_R(x)\defeq\int_{-R}^{R}f(x,t)\,dt$.\\
Since $H_R\to H$ a.e. and $\abs{H_R}\leq h$, dominated convergence gives $\int_XH_R\,d\mu\to\int_XH\,d\mu$.\\
\hlt{Structural Commentary}: Finite interchanges pass to improper limits under an $L^1$ bound.\\
\hlt{Key Insight}: Truncate in $t$, interchange on finite intervals, then pass $R\to+\infty$ by domination in $x$.
\end{proof}
\end{corollary}

\begin{example} \label{ex:differentiating-under-integral-dominated-convergence} \hlt{Differentiating Under the Integral Sign}\\
For $t>0$, set $F(t)\defeq\int_0^\infty x^2e^{-tx}\,dx$. For $t_0>0$ and $\abs{h}<t_0/2$, put $Q_h(x)\defeq\frac{x^2e^{-(t_0+h)x}-x^2e^{-t_0x}}{h}$\\
Then $Q_h\to -x^3e^{-t_0x}$ pointwise and $\abs{Q_h}\leq x^3e^{-(t_0/2)x}\in\mathscr{L}^1(0,\infty)$. Hence $F'(t_0)=-\int_0^\infty x^3e^{-t_0x}\,dx$.\\
\hlt{Structural Commentary}: Parameter differentiation is a limit-interchange problem.\\
\hlt{Key Insight}: Prove a local-in-parameter integrable majorant for the difference quotients.
\end{example}

\begin{example} \label{ex:laplace-transform-powers-factorials} \hlt{Laplace Transform Powers and Factorials}\\
For $t>0$, set $I(t)\defeq\int_0^\infty e^{-tx}\,dx$. Then $I(t)=\lim_{b\to+\infty}\frac{1-e^{-tb}}{t}=t^{-1}$.\\
For $[a,b]\subset(0,+\infty)$ and $k\geq0$, $0\leq x^ke^{-tx}\leq x^ke^{-ax}\in\mathscr{L}^1(0,\infty)$, so repeated differentiation under the integral gives $I^{(n)}(t)=(-1)^n\int_0^\infty x^ne^{-tx}\,dx$. Since $I^{(n)}(t)=(-1)^nn!t^{-n-1}$, for every $n\in\N\cup\{0\}$, $\int_0^\infty x^ne^{-x}\,dx=n!$.\\
\hlt{Structural Commentary}: A parameter integral turns powers of $x$ into derivatives in $t$.\\
\hlt{Key Insight}: The exponential envelope $e^{-ax}$ controls all nearby parameter values $t\geq a$.
\end{example}

\begin{proposition} \label{prop:uniform-convergence-finite-measure-integrals} \hlt{Uniform Convergence on Finite Measure Spaces}\\
Assume $\mu(X)<+\infty$.
\begin{enumerate}[label=(\roman*)]
\setlength{\itemsep}{0pt}
\item If $0\leq f_n$ are measurable and $f_n\rightrightarrows f$, then $f\geq0$ is measurable and $\int_X f\,d\mu=\lim_n\int_X f_n\,d\mu$ in $[0,+\infty]$.
\item If $f_n\in\mathscr{L}^1(\mu)$ and $f_n\rightrightarrows f$, then $f\in\mathscr{L}^1(\mu)$ and $\int_Xf_n\,d\mu\to\int_Xf\,d\mu$ in $\R$.
\end{enumerate}
\begin{proof}
\begin{enumerate}[label=(\roman*)]
\setlength{\itemsep}{0pt}
\item Measurability and $f\geq0$ pass to pointwise limits. For $\varepsilon>0$ and $n\gg1$, $f\leq f_n+\varepsilon\mathbf{1}_X$ and $f_n\leq f+\varepsilon\mathbf{1}_X$.\\
Thus $\int_X f\,d\mu\leq \int_X f_n\,d\mu+\varepsilon\mu(X)$ and $\int_X f_n\,d\mu\leq \int_X f\,d\mu+\varepsilon\mu(X)$. Let $\varepsilon\downarrow0$.
\item Choose $N$ with $\abs{f-f_N}\leq1$. Then $f=f_N+(f-f_N)\in\mathscr{L}^1(\mu)$, and\\
$\abs{\int_Xf_n\,d\mu-\int_Xf\,d\mu}\leq\int_X\abs{f_n-f}\,d\mu\leq\mu(X)\sup_X\abs{f_n-f}\to0$.
\end{enumerate}
\hlt{Structural Commentary}: On finite measure spaces, uniform error controls integral error.\\
\hlt{Key Insight}: Finiteness of $\mu(X)$ turns a supremum bound into an integral bound.
\end{proof}
\end{proposition}

\begin{example} \label{ex:uniform-convergence-needs-finite-measure} \hlt{Why the Finite-Measure Hypothesis Matters}\\
On $(\R,m)$, let $f_n\defeq n^{-1}\mathbf{1}_{[0,n]}$. Then $\norm{f_n}_\infty=1/n\to0$, but $\int_\R f_n\,dm=1$ $\forall n$.\\
Thus uniform convergence need not commute with integrals when $m(\R)=+\infty$; Fatou only yields $0\leq1$.
\end{example}

\begin{example} \label{ex:dominated-convergence-needs-domination} \hlt{Why Domination Cannot Be Dropped}\\
On $(\R,m)$, let $f_n\defeq n\mathbf{1}_{[1/n,2/n]}$. Then $f_n\to0$ pointwise, but $\int_\R f_n\,dm=1$ $\forall n$.\\
Thus pointwise convergence without domination need not commute with integrals; Fatou only yields $0\leq1$.
\end{example}

\begin{theorem} \label{thm:l1-absolutely-summable-series-integrate-termwise} \hlt{Absolutely Summable $L^1$ Series}\\
Let $f_n\in\mathscr{L}^1(\mu)$ and $\sum_{n\geq1}\norm{f_n}_1<+\infty$.\\
Then $\sum_nf_n$ converges absolutely a.e. to some $f\in\mathscr{L}^1(\mu)$, and $\int_X f\,d\mu=\sum_{n\geq1}\int_X f_n\,d\mu$.
\begin{proof}
Set $G\defeq\sum_{n\geq1}\abs{f_n}$. Then $\int_XG\,d\mu=\sum_{n\geq1}\norm{f_n}_1<+\infty$, so $G<+\infty$ a.e. and $G\in\mathscr{L}^1$.\\
Let $f\defeq\sum_{n\geq1}f_n$ a.e. and $S_N\defeq\sum_{n=1}^{N}f_n$.\\
Since $\abs{S_N}\leq G$ and $S_N\to f$ a.e., dominated convergence gives $\int_Xf\,d\mu=\lim_N\int_XS_N\,d\mu=\sum_{n\geq1}\int_Xf_n\,d\mu$.\\
\hlt{Structural Commentary}: Absolute summability in $L^1$ gives a.e. absolute convergence and termwise integration.\\
\hlt{Key Insight}: First integrate the nonnegative majorant $\sum\abs{f_n}$.
\end{proof}
\end{theorem}

\begin{proposition} \label{prop:l1-convergence-controls-absolute-integrals} \hlt{$L^1$ Convergence Controls Absolute Integrals}\\
If $f_n,f\in\mathscr{L}^1(\mu)$ and $\norm{f_n-f}_1\to0$, then $\norm{f_n}_1\to\norm{f}_1$; $\int_X\abs{f}\,d\mu=\lim_n\int_X\abs{f_n}\,d\mu$.
\begin{proof}
The reverse triangle inequality gives $\abs{\abs{f_n}-\abs{f}}\leq\abs{f_n-f}$; hence $\abs{\norm{f_n}_1-\norm{f}_1}\leq\norm{f_n-f}_1\to0$.\\
\hlt{Structural Commentary}: $L^1$ convergence is strong enough to converge the $L^1$ sizes.\\
\hlt{Key Insight}: The reverse triangle inequality converts function convergence into norm convergence.
\end{proof}
\end{proposition}

\begin{proposition} \label{prop:finite-positive-integral-localization} \hlt{Finite Integral Localizes to Finite Measure}\\
If $f:X\to[0,+\infty]$ is measurable, $\int_X f\,d\mu<+\infty$, then $\forall\varepsilon>0$ $\exists$ $E\in\mathcal{M}:\quad \mu(E)<+\infty$; $\int_{X\setminus E} f\,d\mu<\varepsilon$.
\begin{proof}
Set $E_n\defeq\{1/n<f\leq n\}$. Then $\mu(E_n)\leq n\int_X f\,d\mu<+\infty$, and $f\mathbf{1}_{E_n}\uparrow f$ a.e.\\
Hence $\int_{E_n}f\,d\mu\uparrow\int_X f\,d\mu$ by monotone convergence, so $\int_{X\setminus E_n}f\,d\mu\downarrow0$; choose $E=E_n$.\\
\hlt{Structural Commentary}: Finite integral mass can be captured on a finite-measure part of the space.\\
\hlt{Key Insight}: Keep the levels where $f$ is bounded above and away from zero.
\end{proof}
\end{proposition}

\begin{proposition} \label{prop:setwise-convergence-from-total-mass-convergence} \hlt{Setwise Convergence from Total Mass Convergence}\\
Let $f_n,f:X\to[0,+\infty]$ be measurable, $f_n\to f$ pointwise, and $\int_X f\,d\mu=\lim\limits_{n\to\infty}\int_X f_n\,d\mu<+\infty$.\\
Then $\forall E\in\mathcal{M}$, $\int_E f\,d\mu=\lim\limits_{n\to\infty}\int_E f_n\,d\mu$.
\begin{proof}
Fix $E\in\mathcal{M}$ and set $E^c\defeq X\setminus E$. By Fatou, $\int_Ef\,d\mu\leq\liminf\limits_n\int_Ef_n\,d\mu$ and $\int_{E^c}f\,d\mu\leq\liminf\limits_n\int_{E^c}f_n\,d\mu$.\\
Since $\int_X f_n\,d\mu=\int_E f_n\,d\mu+\int_{E^c}f_n\,d\mu\to\int_Xf\,d\mu<+\infty$, $\limsup\limits_n\int_Ef_n\,d\mu\leq\int_Xf\,d\mu-\int_{E^c}f\,d\mu=\int_Ef\,d\mu$.\\
Thus $\liminf_n\int_Ef_n\,d\mu=\limsup_n\int_Ef_n\,d\mu=\int_Ef\,d\mu$.\\
\hlt{Structural Commentary}: Total mass plus Fatou on complements fixes each part.\\
\hlt{Key Insight}: A lower bound on the complement is an upper bound on the chosen set.
\end{proof}
\end{proposition}

\begin{definition} \label{def:function-truncates} \hlt{Function Truncates}\\
For $f:X\to\R$ and $n\in\N$, define the \hlt{$n$th truncate}
\begin{align}
T_nf(x)\defeq (-n)\vee(f(x)\wedge n)=
\begin{cases}
-n,&f(x)<-n,\\
f(x),&-n\leq f(x)\leq n,\\
n,&f(x)>n.
\end{cases} \nonumber
\end{align}
\end{definition}

\begin{proposition} \label{prop:truncation-characterization-integrability} \hlt{Truncates Characterize Integrability}\\
Let $f:X\to\R$ be measurable and $f_n\defeq T_nf$. Then $f\in\mathscr{L}^1(\mu)\Longleftrightarrow \sup_n\norm{f_n}_1<+\infty$.\\
Moreover, if $f\in\mathscr{L}^1(\mu)$, then $\int_Xf_n\,d\mu\to\int_Xf\,d\mu$.
\begin{proof}
The map $T_n:\R\to[-n,n]$ is continuous, so $f_n$ is measurable; also $f_n\to f$, $\abs{f_n}\leq\abs{f}$, and $\abs{f_n}\uparrow\abs{f}$.\\
If $f\in\mathscr{L}^1(\mu)$, dominated convergence gives $\int_Xf_n\,d\mu\to\int_Xf\,d\mu$ and $\sup_n\norm{f_n}_1\leq\norm{f}_1$.\\
Conversely, monotone convergence gives $\norm{f}_1=\lim\limits_{n\to\infty}\norm{f_n}_1=\sup\limits_n\norm{f_n}_1<+\infty$, so $f\in\mathscr{L}^1(\mu)$.\\
\hlt{Structural Commentary}: Truncation reduces $L^1$ questions to bounded approximants.\\
\hlt{Key Insight}: Dominated convergence gives limits; monotone convergence reads integrability from $\abs{f_n}\uparrow\abs{f}$.
\end{proof}
\end{proposition}

\begin{example} \label{ex:setwise-convergence-needs-finite-total-mass} \hlt{Why the Finite Total Mass Assumption Matters}\\
Let $X=\N\times\{0,1\}$ with counting measure, set $E\defeq\N\times\{1\}$.\\
Define $f\defeq\mathbf{1}_{\N\times\{0\}}$; $f_n\defeq\mathbf{1}_{\{1,\ldots,n\}\times\{0\}}+n\mathbf{1}_{\{n\}\times\{1\}}$.\\
Then $f_n\to f$ pointwise and $\int_Xf\,d\mu=\lim_n\int_Xf_n\,d\mu=+\infty$; $\int_Ef\,d\mu=0,\qquad \int_Ef_n\,d\mu=n$.\\
Thus Proposition \ref{prop:setwise-convergence-from-total-mass-convergence} can fail when the common total mass is infinite.
\end{example}
